% from ega4-20-fr.tex
\oldpage[IV-4]{245}
isomorphisme.
Soit $v:\operatorname{dom}_S(\omega)\to\Gamma_\omega$ le $S$-morphisme réciproque de la restriction $p^{-1}(\operatorname{dom}_S(\omega))\to\operatorname{dom}_S(\omega)$ de $p$ ; on note parfois $p^{-1}$ le pseudo-$S$-morphisme de $X$ dans $\Gamma_\omega$, classe de $v$ (dont l'application rationnelle associée \sref{IV.20.2.13}[(ii)] est alors \emph{birationnelle}) ; comme $p^{-1}(\operatorname{dom}_S(\omega))$ est le graphe du $S$-morphisme $u=q\circ v:\operatorname{dom}_S(\omega)\to Y$, il est schématiquement dense dans $\Gamma_\omega$ \sref{IV.11.10.3}[(iv)], donc $\omega$ peut être considéré comme le \emph{composé} $q\circ p^{-1}$ \sref{IV.20.3.2}.
Pour toute partie $M$ de l'espace sous-jacent à $X$, on pose parfois $\omega(M)=q(p^{-1}(M))$, ce qui revient alors à considérer $\omega$ comme une application de $X$ dans $\mathfrak{P}(Y)$ (ou, comme disent certains auteurs, une « fonction multiforme »).
On notera que lorsque $x\in\operatorname{dom}_S(\omega)$, $\omega(\{x\})$ est l'ensemble $\{u(x)\}$ ; en général, pour un $x\in X$ quelconque, si l'on désigne par $s$ l'image de $x$ dans $S$, par $Y_s$ la fibre au point $s$ du morphisme structural $Y\to S$, l'ensemble $\omega(\{x\})$ est une partie du préschéma $Y_s$.
\end{env}

\begin{env}[20.4.3]
\label{IV.20.4.3}
Dans tout le reste de ce numéro, nous allons nous borner au cas où $X$ est \emph{réduit}, de sorte qu'il y a alors identité entre pseudo-$S$-morphismes et $S$-\emph{applications rationnelles} \sref{IV.20.2.7}.
En outre, le \emph{graphe} de toute $S$-application rationnelle de $X$ dans $Y$ est alors défini \sref{IV.20.4.1}.
\end{env}

\begin{proposition}[20.4.4]
\label{IV.20.4.4}
Soient $X$ un $S$-préschéma localement intègre, $Y$ un $S$-préschéma localement de type fini et séparé sur $S$, $\omega$ une $S$-application rationnelle de $X$ dans $Y$, $p:\Gamma_\omega\to X$ la projection canonique.
Alors, si $x\in X$ est un point normal tel que l'ensemble $p^{-1}(x)$ contienne un point isolé, $\omega$ est définie au point $x$.
\end{proposition}

\begin{proof}
En effet, la première projection $p_1:X\times_S Y\to X$ est alors un morphisme séparé et localement de type fini, donc il en est de même de sa restriction $p:\Gamma_\omega\to X$, qui est en outre \emph{birationnel} et comme $\Gamma_\omega$ est réduit et $X$ intègre, $\Gamma_\omega$ est \emph{intègre} ; il résulte donc de (\textbf{Err}\textsubscript{\textbf{IV}}, 30) que l'hypothèse sur $x$ entraîne l'existence d'un voisinage ouvert $U$ de $x$ tel que la restriction $p^{-1}(U)\to U$ de $p$ soit un isomorphisme ; d'où la conclusion \sref{IV.20.4.2}.

On notera que l'énoncé \sref{IV.20.4.4} est la formulation originelle donnée par Zariski à son « Main theorem » (à cela près qu'il se restreignait aux schémas algébriques sur un corps de base).
\end{proof}

\begin{proposition}[20.4.5]
\label{IV.20.4.5}
Les hypothèses et notations étant celles de \sref{IV.20.4.4}, supposons en outre $X$ normal, et soient $X'$ un préschéma réduit, $f:X'\to X$ un morphisme localement de type fini et universellement ouvert.
Alors $\omega\circ f$ est défini, et l'on a $\operatorname{dom}_S(\omega\circ f)=f^{-1}(\operatorname{dom}_S(\omega))$ (en d'autres termes, si $x'\in X'$ et $x=f(x')$, alors, pour que $\omega$ soit définie au point $x$, il faut et il suffit que $\omega\circ f$ le soit au point $x'$).
\end{proposition}

\begin{proof}
Comme $X'$ est réduit et $f^{-1}(\operatorname{dom}_S(\omega))$ partout dense dans $X'$ en vertu de \sref{IV.2.4.11}, le composé $\omega\circ f$ est défini ; pour prouver que, lorsque $\omega\circ f$ est définie au point $x'$, $\omega$ l'est au point $x$, on peut évidemment remplacer $X'$ par un voisinage ouvert de $x'$, donc supposer $\omega\circ f$ \emph{partout définie} ; en outre, comme $f(X')$ est ouvert dans $X$, on peut supposer $f$ \emph{surjectif}.
En vertu de l'hypothèse sur $f$, le morphisme $f_{(Y)}:X'\times_S Y\to X\times_S Y$ est alors ouvert, donc $\overline{f_{(Y)}^{-1}(M)}=f_{(Y)}^{-1}(\overline{M})$ pour toute partie $M$ de $X\times_S Y$ ; prenant pour $M$ l'ensemble $\Gamma_u$, où $u:\operatorname{dom}_S(\omega)\to Y$ est la restriction de $\omega$ à $\operatorname{dom}_S(\omega)$, il résulte de la relation précédente et de \sref[I]{I.3.3.12} que l'ensemble sous-jacent à $\Gamma_{\omega\circ f}$ est égal à $f_{(Y)}^{-1}(\Gamma_\omega)$ ; comme on sait
